\section{Hierarchies and Orders}\label{sec:hierarchies}

Because the higher angels enlighten the lower and never the reverse
(\S\ref{sec:q106}), the angels form an ordered multitude, and Aquinas
turns next to their ordering by hierarchies and orders (\latin{de
ordinatione Angelorum secundum hierarchias et ordines}), first of the good (q.~108) and then of the bad (q.~109)
(\ST{I q.~108 pr.}). The word \gk{hierarchia} is the Areopagite's, and
Aquinas explains it from the Greek: \latin{hierarchia dicitur quasi sacer
principatus a hieron, quod est sacrum, et archon, quod est princeps}
(\emph{In II Sent.} d.~9 q.~1 a.~1). Dionysius defines it as ``a sacred
order and science and operation, assimilated, as far as attainable, to the
likeness of God, and conducted to the illuminations granted to it from God,
according to capacity, with a view to the Divine imitation'' (\emph{CH}
3.1); and Aquinas reads its terms as the parts of every sacred rule,
whose end is likeness to God: order expresses the grade of power, science directs, action leads to
the end, and the likeness of God is the end intended (\emph{In II Sent.}
d.~9 q.~1 a.~1).

The Church learned the nine names from Scripture and the three triads from
Dionysius (\S\ref{sec:ch}; the nine orders one by one are set out in
Appendix~\ref{app:orders}). Gregory the Great preached them to the people
of Rome, gathering the nine from the Prophets and from Paul (\emph{Hom.\ in
Ev.} 34.7), and gave each its ministry; Isidore handed Gregory's
account on (\emph{Etymologiae} VII.5); the Master of the \emph{Sentences}
fixed the doctrine for the schools (\emph{Sent.} II d.~9). Aquinas gives
the reasons of the whole. The tradition carries at the same time a holy
reserve about what is hidden. Augustine, asked what the differences of rank
may be among the thrones, dominions, principalities, and powers of the
Apostle, answered: ``let those who are able answer these questions, if they
can also prove their answers to be true; but as for me, I confess my
ignorance'' (\emph{Enchiridion} 58). Dionysius opens his account of the
orders by confessing that how many and of what sort they are ``the deifying
Author of their consecration alone distinctly knows,'' and that we know them
only ``so far as the Godhead has revealed to us, through them'' (\emph{CH}
6). Aquinas makes that confession the principle of a.~3 below: ``our
knowledge of the angels is imperfect.'' Bernard, contemplating the same
names, speaks throughout in the mood of reverent conjecture: ``Let us
suppose~\dots'' (\emph{De consideratione} V.4.8).

\subsection{The Angelic Degrees of Hierarchies and Orders (\ST{I
q.~108})}\label{sec:q108}

The eight articles ask whether all the angels belong to one hierarchy;
whether in one hierarchy there is only one order; whether in one order
there are many angels; whether the distinction of hierarchies and orders is
from nature; what are the names and properties of each order; how the
orders compare with one another; whether the orders will outlast the Day of
Judgment; and whether men are taken up into the angelic orders (\ST{I
q.~108 pr.}).

\paragraph{Three angelic hierarchies (a.~1).} The \emph{sed contra} is
that ``Dionysius (Coel. Hier. vi) distinguishes three hierarchies of
angels.'' Aquinas defines the thing before he divides it. ``Hierarchy means
a `sacred' principality,'' and principality includes two things, ``the
prince himself and the multitude ordered under the prince.'' On the side of
the prince there is but one hierarchy, ``because there is one God, the
Prince not only of all the angels but also of men and all creatures,'' and
so one hierarchy of all rational creatures able to share in holy things,
as Augustine teaches that there are ``two cities, that is, two societies,
one of the good angels and men, the other of the wicked'' (\ST{I q.~108
a.~1 co.}). Augustine's own words are that it is ``not incongruous and
unsuitable to speak of a society composed of angels and men together; so
that there are not four cities or societies,---two, namely, of angels, and
as many of men,---but rather two in all, one composed of the good, the other
of the wicked, angels or men indifferently'' (\emph{De civitate Dei}
XII.1).

On the side of the multitude, a principality is one in so far as its
members can be governed in one and the same way; ``thus, under one king
there are different cities, which are governed by different laws and
administrators.'' Men and angels receive the divine enlightenments in
different ways, ``for the angels receive them in their intelligible purity,
whereas men receive them under sensible signs, as Dionysius says (Coel.
Hier. i).'' Dionysius's sentence is that ``it is not possible that the
supremely Divine Ray should otherwise illuminate us, except so far as it is
enveloped, for the purpose of instruction, in variegated sacred veils''
(\emph{CH} 1.2). There must therefore be a human hierarchy distinct from
the angelic, and in the same manner three angelic hierarchies are
distinguished. The higher angels have a more universal knowledge
(\S\ref{sec:q55}), and the types of things about which the angels are
enlightened can be considered in three ways: ``first as preceding from God
as the first universal principle, which mode of knowledge belongs to the
first hierarchy, connected immediately with God, and, `as it were, placed in
the vestibule of God,' as Dionysius says''; secondly, ``forasmuch as these
types depend on the universal created causes which in some way are already
multiplied,'' which belongs to the second; thirdly, ``forasmuch as these
types are applied to particular things as depending on their causes,'' which
belongs to the lowest (\ST{I q.~108 a.~1 co.}). Dionysius's image is of
the first beings ``marshalled, as it were, in Its very vestibule''
(\emph{CH} 7.2).

The commentary on the \emph{Sentences} gives the same division an image from
a royal court. Of the ministers of one king, some work about the king's
person, \latin{quasi cubicularii et consiliarii et assessores}, and these
are like the first hierarchy; some have offices for the government of the
realm in common, not assigned to this or that province, \latin{ut domini
regalis curiae, et principes militiae, et judices curiae}, and these are
like the second; some are set over the government of a part of the realm,
\latin{ut praepositi et balivi, et hujusmodi officiales minores}, and these
are like the third (\emph{In II Sent.} d.~9 q.~1 a.~3). The three
hierarchies, three orders in each, reflect the Trinity. The Master had
written that there are three threes \latin{ut Trinitatis similitudo in eis
insinuetur impressa} (Peter Lombard, \emph{Sent.} II d.~9 c.~1), and
Aquinas repeats it: \latin{ut sic in omnibus Trinitatis similitudo
reluceat, secundum principium, medium et finem} (\emph{In II Sent.} d.~9
q.~1 a.~3 ad~3).

From the definition follows a correction: ``those err and speak against the
opinion of Dionysius who place a hierarchy in the Divine Persons, and call it
the `supercelestial' hierarchy.'' Among the Persons there is a natural
order, but no hierarchical order, for as Dionysius says, ``The hierarchical
order is so directed that some be cleansed, enlightened, and perfected; and
that others cleanse, enlighten, and perfect''; ``which far be it from us to
apply to the Divine Persons'' (\ST{I q.~108 a.~1 co.}; \emph{CH} 3.2). As
regards the knowledge of God himself, ``Whom all see in one way---that is,
in His essence---there is no hierarchical distinction among the angels''
(ad~2). And all men are one hierarchy because ``all men are of one species,
and have one connatural mode of understanding; which is not the case in the
angels'' (ad~3; \S\ref{sec:q50}).

\paragraph{Three orders in each hierarchy (a.~2).} The \emph{sed contra} is
the Apostle, who says that God set the Man Christ ``above all principality
and power and virtue and dominion'' (Eph 1:21), ``which are the various
orders of the angels, and some of them belong to one hierarchy.'' A
multitude ``would not be ordered, but confused, if there were not in it
different orders. So the nature of a hierarchy requires diversity of
orders.'' The diversity arises from diversity of offices and actions, as in a
city there is one order of judges, another of soldiers, another of those who
work the fields. Yet all may be reduced to three, ``when we consider that
every multitude has a beginning, a middle, and an end''; so in every city
some are highest, as the nobles (\latin{optimates}), some lowest, as the
common people (\latin{vilis populus}), and some between, \latin{populus
honorabilis}; ``and so in every hierarchy Dionysius places three orders''
(\ST{I q.~108 a.~2 co.}).

The objections sharpen the doctrine of the angelic commonwealth. The Master
had said that in heaven ``nothing is possessed individually,'' so that all
spiritual gifts would be common and there would be no grades. Aquinas
answers: ``All things are possessed in common by the angelic society, some
things, however, being held more excellently by some than by others. Each
gift is more perfectly possessed by the one who can communicate it, than by
the one who cannot communicate it; as the hot thing which can communicate
heat is more perfect that [\emph{sic}] what is unable to give heat. And the more
perfectly anyone can communicate a gift, the higher grade he occupies, as he
is in the more perfect grade of mastership who can teach a higher science''
(ad~2). Grade among the angels is measured by the power to give. And the
orders are not distinguished, as the orders of the Church are, by cleansing,
enlightening, and perfecting, for ``the inferior angel is superior to the
highest man of our hierarchy, according to the words, `He that is the lesser
in the kingdom of heaven, is greater than he'---namely, John the Baptist''
(Matt 11:11); ``hence the lesser angel of the heavenly hierarchy can not only
cleanse, but also enlighten and perfect, and in a higher way than can the
orders of our hierarchy'' (ad~3).

\paragraph{Many angels in one order (a.~3).} The \emph{sed contra} is the
vision of Isaiah: ``The Seraphim cried to one another'' (Isa 6:3); there are
therefore many angels in the one order of the Seraphim. ``Whoever knows
anything perfectly, is able to distinguish its acts, powers, and nature,
down to the minutest details, whereas he who knows a thing in an imperfect
manner can only distinguish it in a general way,'' as one who knows nature
imperfectly places the heavenly bodies in one order, inanimate bodies in
another, plants and animals in others. ``Now our knowledge of the angels is
imperfect, as Dionysius says (Coel. Hier. vi). Hence we can only distinguish
the angelic offices and orders in a general way, so as to place many angels
in one order. But if we knew the offices and distinctions of the angels
perfectly, we should know perfectly that each angel has his own office and
his own order among things, and much more so than any star, though this be
hidden from us'' (\ST{I q.~108 a.~3 co.}). Since each angel is his own
species (\ST{I q.~50 a.~4}; \S\ref{sec:q50}), the angels of one order are
equal only ``in a common similitude,'' not absolutely; ``hence Dionysius
says (Coel. Hier. x) that in one and the same order of angels there are
those who are first, middle, and last'' (ad~1). His text reads: ``we observe
every Hierarchy distributed into first, and middle, and last Powers''
(\emph{CH} 10.2). Two angels at the border of two orders are nearer in
nature than either is to the far members of its own order, ``but less akin
in their fitness for similar offices,'' as two neighbouring points on the
border of white and black are nearer in place and farther in quality
(ad~3).

The commentary on the \emph{Sentences} draws from the same principle a
conclusion about beauty. Since no two angels share a species and species in
the angels differ by more and less of act, \latin{non est accipere duos
Angelos aequales in tota caelesti hierarchia}; and its \emph{sed contra}
adds, \latin{in Angelis est summa pulchritudo post Deum. Sed in diversorum
graduum consonantia ordinata consistit ratio pulchritudinis} (\emph{In II
Sent.} d.~9 q.~1 a.~5). The Master had answered the same question with an
image from the ranks of the saints: there is one order of Apostles and another
of Martyrs, \latin{et tamen in Apostolis alii sunt aliis digniores,
similiter et in Martyribus alii aliis sunt superiores; ita et in ordinibus
Angelorum recte creditur esse} (\emph{Sent.} II d.~9 c.~5). How many angels
stand in each order, Aquinas adds, \latin{ignotus est nobis} (\emph{In II
Sent.} d.~9 q.~1 a.~3).

\paragraph{The distinction of the orders is from grace, disposed by nature
(a.~4).} The \emph{sed contra} is the Master's definition of an order: ``an
angelic order is a multitude of heavenly spirits, who are likened to each
other by some gift of grace, just as they agree also in the participation of
natural gifts.'' The Master's text reads: \latin{Ordo Angelorum dicitur
multitudo caelestium spirituum, qui inter se in aliquo munere gratiae
similantur, sicut et in naturalium datorum participatione conveniunt}
(\emph{Sent.} II d.~9 c.~2). Aquinas determines by the end. ``The order of
government, which is the order of a multitude under authority, is derived
from its end.'' The end of the angels may be taken according to the faculty
of their nature, to know and love God by natural knowledge and love, and in
relation to this end ``the orders of the angels are distinguished by natural
gifts.'' Or it may be taken above their nature, ``in the vision of the Divine
Essence, and in the unchangeable fruition of His goodness; to which end they
can reach only by grace; and hence as regards this end, the orders in the
angels are adequately distinguished by the gifts of grace, but dispositively
by natural gifts'' (\latin{completive quidem secundum dona gratuita,
dispositive autem secundum dona naturalia}), because to the angels grace was
given according to the capacity of their nature (\ST{I q.~62 a.~6};
\S\ref{sec:q62}). ``Hence among men the orders are distinguished according to
the gratuitous gifts only, and not according to natural gifts'' (\ST{I
q.~108 a.~4 co.}).

The commentary on the \emph{Sentences} names the two principles:
\latin{donum gratiae est perfectivum naturae; et ideo distinctio ordinum est
per diversum donum gratuitum sicut per principium formale, et per diversum
donum naturale sicut per principium quasi materiale et dispositivum}; and so
the orders were in a manner distinct from the beginning of creation,
\latin{non tamen secundum ultimum complementum} (\emph{In II Sent.} d.~9
q.~1 a.~7). The Master had put the question whether the orders were
distinct from the first moment. Some of every order fell, which seems to
show that the orders were there from the beginning; yet the fallen never
burned with charity nor were God's seat, and so were never Seraphim or
Thrones. He answers that before the fall of some the orders were not yet
such, since the gifts in which their members agree had not been given; but
the angels were created in grades, \latin{superiores, qui natura magis
subtiles et sapientia amplius perspicaces; inferiores, qui natura minus
subtiles et intelligentia minus perspicaces}, and these invisible
differences only he can weigh who has ordered all things in number, measure,
and weight: \latin{Has autem invisibiles differentias invisibilium solus
ille ponderare potest qui omnia in numero et mensura et pondere disposuit}
(\emph{Sent.} II d.~9 c.~4; Wis 11:21). Aquinas, who judges it more probable
that all the angels were created in grace (\ST{I q.~62 a.~3};
\S\ref{sec:q62}), can say that the orders were in a manner there from the beginning;
the Master, who holds grace added afterward to those who stood, says that
the orders came with the gifts.

\paragraph{The names of the orders (a.~5).} The \emph{sed contra} is ``the
authority of Holy Scripture wherein they are so named'': Seraphim in Isaiah
6:2, Cherubim in Ezekiel, Thrones in Colossians 1:16, Dominations, Virtues,
Powers, and Principalities in Ephesians 1:21, Archangels in the Epistle of
Jude (``Michael the archangel, disputing with the devil,'' Jude 9), and
Angels in many places. Aquinas takes from Dionysius the rule that ``the
proper name of each order expresses its property'' and gives the logic of
such naming. Something is found in another by way of property when it is
adequate to its nature, by way of excess when it is possessed eminently, and
by way of participation when it is possessed partially, ``thus holy men are
called gods by participation.'' A thing is named properly from what is
adequate to it, as man is properly ``a `rational substance,' but not an
`intellectual substance,' which latter is the proper name of an angel;
because simple intelligence belongs to an angel as a property, and to man by
participation.'' So among the angels ``all spiritual perfections are common
to all the angels, and~\dots\ they are all more excellently in the superior
than in the inferior''; and since among these perfections there are grades,
``the superior perfection belongs to the superior order as its property,
whereas it belongs to the inferior by participation; and conversely the
inferior perfection belongs to the inferior order as its property, and to
the superior by way of excess; and thus the superior order is denominated
from the superior perfection'' (\ST{I q.~108 a.~5 co.}).

Two masters expound the names. ``Dionysius (Coel. Hier. vii) explains the
names of the orders accordingly as they befit the spiritual perfections they
signify. Gregory, on the other hand, in expounding these names~\dots\ seems
to regard more the exterior ministrations; for he says that `angels are so
called as announcing the least things; and the archangels in the greatest;
by the virtues miracles are wrought; by the powers hostile powers are
repulsed; and the principalities preside over the good spirits themselves'\,''
(\ST{I q.~108 a.~5 co.}). Gregory begins his exposition with the principle
that has become a commonplace of the Latin tradition, \latin{angelorum
vocabulum nomen est officii, non naturae} (\emph{Hom.\ in Ev.}\ 34.8;
\S\ref{sec:fa-gregory}; Isidore repeats it, \emph{Etym.} VII.5.2).
Aquinas reconciles the two masters on this very word. As naming the
condition of a nature that manifests what is above it, ``angel'' is a name
of nature, and the lowest order keeps it because it manifests all the higher
orders to us; \latin{Sed si nuntiare sumatur secundum exterius ministerium
quod exercent, sic Angelus est nomen officii, non naturae, ut dicit
Gregorius} (\emph{In II Sent.} d.~9 q.~1 a.~4 ad~2).

The replies then take the names one by one, and each reply teaches
something about the structure of the whole. The common name ``angel'' remains
proper to the lowest order, ``as Dionysius says (Coel. Hier. v),'' or because
they ``announce things to us immediately''; and ``virtue'' is common to all
as the medium between essence and operation, but proper to one order as ``a
certain excellence of strength,'' which Dionysius calls ``a certain virile
and immovable strength'' (ad~1). Parker's Dionysius has the theologians, when
they describe the orders more accurately, ``name exclusively, `angelic rank,'
that which completes the full tale of the Divine and Heavenly Hosts''
(\emph{CH} 5), and speak of all the divine minds as ``distributed into
three,---into essence, and power, and energy'' (\emph{CH} 11.2). Dominion
belongs to God ``in a special manner, by way of excess: but the Divine word
gives the more illustrious heavenly princes the name of Lord by
participation, through whom the inferior angels receive the Divine gifts''
(ad~2, citing \emph{DN} 12). Domination, Power, and Principality belong to
government in different ways: the lord only prescribes; power signifies
ordination, ``He that resisteth the power, resisteth the ordination of God''
(Rom 13:2); to preside is ``to be first among others,'' as ``Princes went
before joined with singers'' (Ps 67[68]:26) (ad~3). Gregory's sentence is
\latin{Nam principari est inter reliquos priorem existere: dominari vero est
etiam subjectos quosque possidere} (\emph{Hom.\ in Ev.} 34.10). The
Archangels stand between the Principalities and the Angels, and ``a medium
compared to one extreme seems like the other,'' as tepid seems cold beside
hot and hot beside cold; so they are called ``angel princes'' (ad~4). The
Seraphim are named not from charity alone but ``from the excess of charity,
expressed by the word ardor or fire,'' and in fire Aquinas finds the upward
and continuous movement by which they are borne unswervingly to God, the
penetrating heat by which they rouse those below them ``to a like fervor,
and cleansing them wholly by their heat,'' and the brightness by which ``these
angels have in themselves an inextinguishable light, and~\dots\ also
perfectly enlighten others''; the Cherubim are named from ``a certain excess
of knowledge,'' which Dionysius unfolds in four things, the perfect vision of
God, the full reception of the divine light, the contemplation in God of the
beauty of the divine order, and the copious outpouring of this knowledge on
others (ad~5). The Thrones excel the lower orders ``as having an immediate
knowledge of the types of the Divine works''; they are like material seats
in four ways: raised above the earth, made firm (``for the angels themselves
are made firm by God''), receiving God and in a certain way bearing him to
the lower creatures, and open on one side, ``thus are the angels promptly
open to receive God and to serve Him'' (ad~6). For each of these names
Dionysius's own exposition stands in \S\S\ref{sec:ch7}--\ref{sec:ch9}.

Bernard of Clairvaux gives the names their most beautiful contemplative
reading, in Gregory's order, and turns each name toward God (\emph{De
consideratione} V.4.8--10; expounded at \S\ref{sec:fa-bernard}). His
Dominions tower above the rest, so that to them, ``as it were to their
masters, the Principalities account for their commands, the Powers for their
defences, the Virtues for their operations, the Archangels for their
revelations, the Angels for their care and foresight'' (V.4.8); and he closes
the contemplation with the law of participation that Aquinas states in the
language of the schools: ``The Seraphim burn, but with the fire of God, or rather with fire for God. Their
chief characteristic is their love, but they love not as much, nor in the
same way, as God. The Cherubim shine, and excel in knowledge, but by
participation in the truth~\dots\ Angels and Archangels are with us, but He
is more our own who is not only with us but in us'' (V.5.11).

\paragraph{The ordering of the orders (a.~6).} The \emph{sed contra} is the
Dionysian sequence: in the highest hierarchy the Seraphim first, the
Cherubim in the middle, the Thrones last; in the middle hierarchy the
Dominations first, the Virtues in the middle, the Powers last; in the lowest
the Principalities first, then the Archangels, and last the Angels (\ST{I
q.~108 a.~6 s.c.}; \emph{CH} 7--9). Aquinas states the difference at once.
``The grades of the angelic orders are assigned by Gregory~\dots\ and
Dionysius~\dots, who agree as regards all except the `Principalities' and
`Virtues.'\,'' Dionysius places the Virtues beneath the Dominations and above
the Powers, and the Principalities beneath the Powers and above the
Archangels; Gregory places the Principalities between the Dominations and
the Powers, and the Virtues between the Powers and the Archangels. ``Each of
these placings may claim authority from the words of the Apostle.'' To the
Ephesians Paul numbers the middle orders ascending, ``above all principality
and power and virtue and dominion'' (Eph 1:21), placing the Virtues between
Powers and Dominations, as Dionysius does; to the Colossians he numbers them
descending, ``whether thrones, or dominations, or principalities, or powers''
(Col 1:16), placing the Principalities between Dominations and Powers, as
Gregory does (\ST{I q.~108 a.~6 co.}; the two sequences are set out in
Appendix~\ref{app:orderings}).

Aquinas examines Dionysius's arrangement first. The first hierarchy
contemplates the ideas of things in God himself, the second in the universal
causes, the third in their application to particular effects. ``And because
God is the end not only of the angelic ministrations, but also of the whole
creation, it belongs to the first hierarchy to consider the end; to the
middle one belongs the universal disposition of what is to be done; and to
the last belongs the application of this disposition to the effect, which is
the carrying out of the work.'' Dionysius therefore places in the first
hierarchy the orders named from their relation to God, in the middle those
whose names denote ``a certain kind of common government or disposition,''
and in the last those whose names denote ``the execution of the work.''
Within the first, as God is the end of creatures ``as the leader is the end
of an army,'' Aquinas finds the likeness of a court: ``there are some who
enjoy the dignity of being able with familiarity to approach the king or
leader; others in addition are privileged to know his secrets; and others
above these ever abide with him, in a close union.'' So the Thrones receive
God familiarly in themselves, knowing immediately the types of things in
him, which belongs to the whole first hierarchy; ``the `Cherubim' know the
Divine secrets supereminently; and the `Seraphim' excel in what is the
supreme excellence of all, in being united to God Himself.'' In government
the Dominations appoint what is to be done, the Virtues give the power of
carrying it out, and the Powers order how what has been commanded can be
carried out by others. In execution, which consists in announcing divine
things, there are leaders, ``as in singing, the precentors; and in war,
generals and officers,'' who are the Principalities; those who simply
execute, the Angels; and between them the Archangels (\ST{I q.~108 a.~6
co.}).

The arrangement is ``quite a reasonable one,'' for ``the highest in an
inferior order always has affinity to the lowest in the higher order; as the
lowest animals are near to the plants.'' Here Aquinas lifts the eye to the
highest order of all. ``Now the first order is that of the Divine Persons,
which terminates in the Holy Ghost, Who is Love proceeding, with Whom the
highest order of the first hierarchy has affinity, denominated as it is from
the fire of love.'' The Thrones, lowest of the first hierarchy, are akin to
the Dominations, being so called, according to Gregory, ``because through
them God accomplishes His judgments,'' and enlightened so as to enlighten the second
hierarchy immediately; and the Powers are akin to the Principalities, who
preside ``over the government of peoples and kingdoms,'' the first and
principal of the divine ministrations, ``for the good of a nation is more
divine than the good of one man''; ``and hence it is written, `The prince of
the kingdom of the Persians resisted me' (Daniel 10:13)'' (\ST{I q.~108
a.~6 co.}).

Gregory's arrangement is ``also reasonable.'' Since the Dominations appoint
the divine ministrations, the orders under them are arranged according to
the things in which those ministrations are exercised, and Aquinas takes the
order of those things from Augustine. Augustine's text reads: ``as the more
gross and inferior bodies are governed in due order by the more subtle and
powerful ones, so all bodies are governed by the living spirit; and the
living spirit devoid of reason, by the reasonable living spirit; and the
reasonable living spirit that makes default and sins, by the living and
reasonable spirit that is pious and just; and that by God Himself''
(\emph{De Trinitate} III.4.9). So after the Dominations come the
Principalities, ``who rule even over good spirits; then the `Powers,' who
coerce the evil spirits; even as evil-doers are coerced by earthly powers''
(Rom 13:3--4); then the Virtues, ``which have power over corporeal nature in
the working of miracles''; and then the Archangels and Angels, ``who announce
to men either great things above reason, or small things within the purview
of reason'' (\ST{I q.~108 a.~6 co.}). Gregory's own homily defines the
Principalities as those \latin{qui ipsis quoque bonis angelorum spiritibus
praesunt, qui subjectis aliis dum quaeque sunt agenda disponunt, eis ad
explenda divina ministeria principantur}, and the Powers as those to whose
authority the adverse powers are subject, \latin{quorum potestate
refrenantur, ne corda hominum tantum tentare praevaleant, quantum volunt}
(\emph{Hom.\ in Ev.} 34.10).

The replies settle the ranking of knowledge and love. The orders named from
presiding are not the highest, because ``the angel's subjection to God is
greater than their presiding over inferior things; and the latter is derived
from the former'' (ad~1). The nearness to God signified by the Thrones
belongs to the Cherubim and Seraphim more excellently (ad~2). And though
knowledge seems to come before love, the Seraphim stand above the Cherubim:
``knowledge takes place accordingly as the thing known is in the knower; but
love as the lover is united to the object loved. Now higher things are in a
nobler way in themselves than in lower things; whereas lower things are in
higher things in a nobler way than they are in themselves. Therefore to know
lower things is better than to love them; and to love the higher things, God
above all, is better than to know them'' (ad~3; on the angels' love,
\S\ref{sec:will}). Of the two arrangements Aquinas concludes: ``A careful
comparison will show that little or no difference exists in reality between
the dispositions of the orders according to Dionysius and Gregory.''
Gregory's Principalities, presiding over good spirits, answer to Dionysius's
Virtues, which give efficacy to the lower spirits; and Gregory's Virtues,
working miracles, answer to Dionysius's Principalities, ``for to work
miracles holds the first place in the Divine ministrations; since thereby the
way is prepared for the announcements of the `Archangels' and the `Angels'\,''
(ad~4).

The \emph{Summa contra gentiles} gives the fullest form of the same
architecture. The highest intellects receive the reason of the order of
providence in the last end, which is the divine goodness, and these are the
Seraphim, \latin{quasi ardentes vel incendentes, quia per incendium designari
solet intensio amoris vel desiderii, quae sunt de fine}; the second know it
in the divine form itself, the Cherubim, \latin{quod interpretatur scientiae
plenitudo: scientia enim per formam scibilis perficitur}; the third consider
the disposition of the divine judgments in itself, the Thrones, \latin{nam
per thronum potestas iudiciaria designatur}, as it is written, ``thou hast
sat on the throne, who judgest justice'' (Ps 9:5) (\emph{SCG} III.80). The
middle hierarchy knows the order of providence in universal causes: the
Dominations distribute it to many executors, \latin{dominorum enim est
praecipere quid alii exequantur}; the Virtues multiply it into various
effects, and to them belong the motion of the heavens and the working of
miracles, \latin{nam ista sunt sublimissima in divinis ministeriis}; the
Powers keep the established order unconfused by restraining what could
disturb it. The lowest are set over human affairs: the Principalities over
the common good of cities and nations, so that \latin{dispositio regnorum, et
mutatio dominationis a gente in gentem, ad ministerium huius ordinis
pertinere oportet}; the Archangels over a human good that belongs to one yet
serves many, \latin{sicut ea quae sunt fidei, et cultus divinus}, as Gabriel
announced to the Virgin the Incarnation of the Word, \latin{ab omnibus
credendam}; the Angels over the good of each single man, whence they are called
the guardians of men, ``For he hath given his angels charge over thee''
(Ps 90[91]:11). To those who look carefully, Gregory's arrangement differs
little (\latin{sed, diligenter inspicientibus, utraque ordinatio in modico
differt}), and both have the
authority of the Apostle. And the chapter closes with the law that binds the
whole: \latin{Est autem in omnibus ordinatis virtutibus hoc commune, quod in
vi superioris virtutis omnes inferiores agunt. Unde id quod diximus ad
Seraphim ordinem pertinere, omnes inferiores ex virtute ipsius exequuntur}
(\emph{SCG} III.80). What the Seraphim are, every lower order carries out in
the power of the Seraphim.

The commentary on the \emph{Sentences} had judged more sharply between the
two masters: \latin{ordinatio Dionysii est rationabilior}, since Dionysius
distinguishes the second hierarchy from the third as universal from
particular, unlimited from limited; \latin{Gregorius autem non multum
intendit ordinare, sed numerare}, though his order too can be taken from the
exterior acts and from conformity to our hierarchy, in which those who have
prelacy stand above private persons (\emph{In II Sent.} d.~9 q.~1 a.~3
ad~7). The same article reasons the orders of each hierarchy differently from
the \emph{Summa}. The first hierarchy is numbered by three acts of the
soul's fruition of God: to hold, which perfects memory, and so the Thrones,
in whom God sits and rests; to see, which perfects intellect, and so the
Cherubim; to love, which perfects will, and so the Seraphim, first because
\latin{amor facit interiora amati penetrari}. The second is numbered by the
order of peace in a commonwealth: a sentence determining what is due to
each, which belongs to the Powers; efficacy, so that nothing determined is
impossible to execute, which belongs to the Virtues, as secular laws must be
armed to have coercive force; and rightness under one who directs, which
belongs to the Dominations, \latin{sicut architector in mechanicis}. The
third is numbered by limited power: over a province, the Principalities, as
the prince of the Persians; over one man, the Angels, \latin{quorum est
etiam singulos custodire}; over a good of many fulfilled through one person,
the Archangels, \latin{sicut Gabriel qui nuntiavit nativitatem Christi, et
Joannis Baptistae} (\emph{In II Sent.} d.~9 q.~1 a.~3).

The Latin tradition before Aquinas had carried the two arrangements side by
side. The Master names Dionysius as the author of the three triads and then
lists them in the order Gregory preached: \latin{superiores: Seraphim,
Cherubim, Throni; medii: Dominationes, Principatus, Potestates; inferiores:
Virtutes, Archangeli, Angeli} (\emph{Sent.} II d.~9 c.~1). Bernard, in the
passage above, ascends through Gregory's sequence. Isidore sets the
Dominations above both Virtues and Principalities, \latin{qui etiam
Virtutibus et Principatibus praeeminent} (\emph{Etym.} VII.5.20). Both
orderings thus passed into the tradition side by side; Aquinas judged that
their difference touches the arrangement of the names more than the reality
of the offices.

\paragraph{The orders remain after the Judgment (a.~7).} The \emph{sed
contra} is the song of Deborah, ``the stars, remaining in their order and
courses'' (Judg 5:20), ``which is applied to the angels.'' Two things are
considered in the angelic orders, ``the distinction of grades, and the
execution of their offices.'' The distinction of grades follows nature and
grace, and both remain, ``for these differences of natures cannot be taken
from them unless they themselves be corrupted''; ``the difference of glory
will also ever remain in them according to the difference of preceding
merit.'' The execution of the offices ``will to a certain degree remain after
the Day of Judgment, and to a certain degree will cease. It will cease
accordingly as their offices are directed towards leading others to their
end; but it will remain, accordingly as it agrees with the attainment of the
end. Thus also the various ranks of soldiers have different duties to perform
in battle and in triumph'' (\ST{I q.~108 a.~7 co.}). When the Apostle says
that Christ will ``have brought to nought all principality and power and
virtue'' when he delivers the kingdom to the Father (1~Cor 15:24), he speaks
of their office of leading others to the end, ``because when the end is
attained, it is no longer necessary to tend towards the end'' (ad~1).

The second reply is a small treatise on the dependence of knowledge. In us,
the knowledge of a conclusion depends on the middle terms not only in its
first acquisition ``but also as regards the keeping of the knowledge
acquired,'' for whoever forgets one of the middle terms may keep opinion or
belief about the conclusion, but not science. ``So, since the inferior angels
know the types of the Divine works by the light of the superior angels, their
knowledge depends on the light of the superior angels not only as regards the
acquisition of knowledge, but also as regards the preserving of the knowledge
possessed''; and so after the Judgment, though they will not progress in the
knowledge of some things, ``still this will not prevent their being
enlightened by the superior angels'' (ad~2). Those already saved ``will be
enlightened through the angelic ministry'' (ad~3). The commentary on the
\emph{Sentences} states the same distinction: the essentials of hierarchy
remain for ever, \latin{quia unus alium illuminabit, purgabit, et
perficiet}; the prelacy of the orders by which men are led back to God
through them will cease, \latin{cessante statu viatorum} (\emph{In II Sent.}
d.~11 q.~2 a.~6).

\paragraph{Men are taken up into the angelic orders (a.~8).} The
\emph{sed contra} is the Lord's word of the saints, ``they~\dots\ shall be as
the angels of God in heaven'' (Matt 22:30). Considered only by grade of
nature, ``men can in no way be assumed into the angelic orders; for the
natural distinction will always remain.'' Some, considering this, denied that
men could ever be made equal to the angels; ``but this is erroneous,
contradicting as it does the promise of Christ saying that the children of
the resurrection will be equal to the angels in heaven (Luke 20:36).'' In
Douay's words the children of the resurrection ``are equal to the angels and
are the children of God.'' ``For whatever belongs to nature is the material
part of an order; whilst that which perfects is from grace which depends on
the liberality of God, and not on the order of nature. Therefore by the gift
of grace men can merit glory in such a degree as to be equal to the angels,
in each of the angelic grades; and this implies that men are taken up into
the orders of the angels.'' Others say that not all the saved are taken up
into the orders, ``but only virgins or the perfect; and that the other will
constitute their own order, as it were, corresponding to the whole society of
the angels''; this is against Augustine, who teaches ``that `there will not
be two societies of men and angels, but only one; because the beatitude of
all is to cleave to God alone'\,'' (\ST{I q.~108 a.~8 co.}). Augustine's page
reads: ``those who have this good in common, have, both with Him to whom they
draw near, and with one another, a holy fellowship, and form one city of
God---His living sacrifice, and His living temple,'' of which a part ``is
hereafter to be united to the immortal angels, and which at present is being
gathered from among mortal men'' (\emph{De civitate Dei} XII.9).

The replies mark the limits. Grace is given to the angels in proportion to
their nature, but not so to men; ``so, as the inferior angels cannot be
transferred to the natural grade of the superior, neither can they be
transferred to the superior grade of grace; whereas men can ascend to the
grade of grace, but not of nature'' (ad~1). The angels, by the order of
nature, stand between us and God, and by common law administer human affairs
and all bodies; holy men after this life remain of our nature, and by common
law ``do not administer human affairs, `nor do they interfere in the things
of the living,' as Augustine says''; yet ``by a certain special dispensation
it is sometimes granted to some of the saints to exercise these offices; by
working miracles, by coercing the demons, or by doing something of that
kind'' (ad~2). Augustine had concluded from the miracles at the shrines of
the martyrs that ``it must needs be by a Divine power that the Martyrs are
interested in affairs of the living, from the very fact that for the departed
to be by their proper nature interested in affairs of the living is
impossible'' (\emph{De cura pro mortuis gerenda} 16.19). Finally, it is not
erroneous to say that wicked men are transferred to the penalty of the demons;
the error is that ``the demons are nothing but souls of the dead; and it is
this that Chrysostom rejects'' (ad~3). Chrysostom, preaching on the demoniacs
among the tombs, says that the demons ``would fain suggest to the multitude a
pernicious opinion, as though the souls of the dead become demons, which God
forbid we should ever admit into our conception'' (\emph{Hom.\ in Matt.}
28.3).

Gregory's homily on the ten drachmas is the great patristic text of this
article. The woman had ten drachmas because there are nine orders of angels,
and man was created the tenth, \latin{ut compleretur electorum numerus}
(\emph{Hom.\ in Ev.}\ 34.6); and since the heavenly city is made of angels
and men, men are assigned to the several choirs by the likeness of their
lives, \latin{in eorum sortem per conversationis similitudinem deputantur},
from those who faithfully announce little things to their brethren, who run
into the number of the Angels, to those who burn with contemplation and
kindle others, whose heart, \latin{in igne conversum}, gives light and burns
as the Seraphim's (34.11--12; the whole ladder is expounded at
\S\ref{sec:fa-gregory}).

The Master receives Gregory's teaching and rejects the reading by which the
lost drachma would make a tenth order: \latin{non esse de hominibus
formandum decimum ordinem, tanquam novem sint Angelorum et decimus hominum,
sed homines pro qualitate meritorum statuendos in ordinibus Angelorum}; the
tenth order is said \latin{compleri ex hominibus} in the sense that what fell
among the angels is restored from men (\emph{Sent.} II d.~9 c.~6). Aquinas's commentary on
the passage sets out three positions. The Greek commentator Eustratius held
that no soul can hope to reach the operation of the separate intellects, so
that the glorified would form a tenth order beneath the nine; \latin{Sed haec
positio est contraria dictis sanctorum, et videtur sapere haeresim, cum
etiam beata virgo super choros Angelorum exaltata sit}, unless it be taken of
what man can reach by his natural powers. A second position assigns virgins,
or the perfect, to the angelic orders and the rest to a tenth; but many who
were not virgins, \latin{ut Petrus et Magdalena}, surpass many virgins, and
since from men and angels there is to be one Church and one hierarchy, the
number of orders, three by three as a trace of the uncreated Trinity, will
not be increased by men. \latin{Unde tertia positio plus mihi placet, quae
etiam dictis sanctorum magis consonat, scilicet quod omnes electi assumantur
ad ordines Angelorum, quidam ad superiores, quidam ad inferiores, quidam ad
medios pro diversitate suorum meritorum; sed beata virgo Maria super omnes}
(\emph{In II Sent.} d.~9 q.~1 a.~8). Man's merit can raise him so high, the
commentary adds, because it is in a way more efficacious than the angel's:
for its difficulty, for the continual increase of grace in a longer way of
pilgrimage, and because our merits draw their efficacy from the merit of
Christ, \latin{nusquam enim Angelos apprehendit, sed semen Abrahae} (ad~2;
Heb 2:16). The Old Law's hierarchy was the way into the Church's and was
taken up into it; so ours is the way into the heavenly, \latin{unde in patria
non erit alia hierarchia hominum et Angelorum, sed una et eadem et homines in
ordines Angelorum distribuentur} (ad~4). On the Blessed Virgin above the
choirs, see \S\ref{sec:christ-mary}.

On the number of men to be taken up the Doctors hold two positions. Gregory teaches
that as many men will ascend as there remained elect angels: \latin{ad quam
tantum credimus humanum genus ascendere, quantos illic contigit electos
angelos remansisse} (\emph{Hom.\ in Ev.} 34.11), and the Master adopts
it, \latin{Non enim iuxta numerum eorum qui ceciderunt, sed eorum qui
permanserunt, homines ad beatitudinem admittuntur} (\emph{Sent.} II d.~9
c.~7). Others hold that men replace the number of the fallen, and the Master
reads Augustine that way. Augustine teaches that the restored part of
mankind ``should fill up the gap which the rebellion and fall of the devils
had left in the company of the angels,'' so that the city of God ``shall not
be spoiled of any of the number of her citizens, shall perhaps reign over
even a more abundant population''; ``We do not know the number either of the
saints or of the devils'' (\emph{Enchiridion} 29). Aquinas refers the
number to God, as Augustine does: whether as many men are taken up as fell,
or as stood, or both, or more, or fewer, \latin{ille scit cui soli cognitus
est numerus electorum, in superna felicitate locandus} (\emph{In II Sent.}
d.~9 q.~1 a.~8).

\angeldossier{\ST{I q.~108 aa.~1--8}; parallels \emph{In II Sent.} d.~9
q.~1 aa.~1--8, d.~11 q.~2 a.~6; \emph{Summa contra gentiles} III.80.}
 {Common teaching that the angels are nine orders in three hierarchies
 (aa.~1--2; the names from Scripture, Isa 6:2, Ezek 10, Col 1:16, Eph 1:21,
 Jude 9; the triads from Dionysius, \emph{CH} 6; Gregory, \emph{Hom.\ in Ev.}\ 34.7;
 Peter Lombard, \emph{Sent.} II d.~9 c.~1). Thomist position on the ground
 of the three hierarchies in three grades of universal knowledge (a.~1) and
 on the reduction of each hierarchy to beginning, middle, and end (a.~2).
 Common teaching that there are many angels in each order, with first,
 middle, and last (a.~3; \emph{CH} 10.2); Thomist position that, known
 perfectly, each angel has his own order (a.~3; \ST{I q.~50 a.~4}). Thomist
 position that the orders are distinguished completively by grace and
 dispositively by nature (a.~4); the Master teaches that the orders came
 with the gifts given to those who stood (\emph{Sent.} II d.~9 c.~4). Common
 teaching that the names signify properties of the orders (a.~5; Dionysius,
 \emph{CH} 7--9; Gregory, \emph{Hom.\ in Ev.}\ 34.8--10). Disputed: the ordering of
 the middle orders, Dionysius placing the Virtues second and the
 Principalities seventh, Gregory the Principalities fifth and the Virtues
 seventh (a.~6), with Aquinas's judgment that they differ little in reality
 (a.~6 ad~4; \emph{SCG} III.80) and his earlier preference for Dionysius
 (\emph{In II Sent.} d.~9 q.~1 a.~3 ad~7). Common teaching that the orders
 remain for ever in glory (a.~7); Thomist position on the offices that
 cease and remain. Scripture teaches that the children of the resurrection
 are equal to the angels (Luke 20:36). Common teaching that the elect are
 taken up into the angelic orders according to their merits, not into a
 tenth order (a.~8; Gregory, \emph{Hom.\ in Ev.}\ 34.11; Lombard, \emph{Sent.} II
 d.~9 c.~6; \emph{In II Sent.} d.~9 q.~1 a.~8), and that the Blessed Virgin
 is exalted above all the choirs (\emph{In II Sent.} d.~9 q.~1 a.~8).
 Disputed whether the elect of men equal the number of angels who stood or
 of those who fell (Gregory and the Master; Augustine, \emph{Enchiridion}
 29, as the Master reads him); the number is known to God alone.}
 {Dionysius, \emph{CH} 6 (\emph{sed contra} of a.~1; a.~3 co.), \emph{CH} 7
 (a.~1 co., a.~5 co., ad~5--6, \emph{sed contra} of a.~6), \emph{CH} 1 and 3
 (a.~1 co.), \emph{CH} 8 and 9 (a.~5 ad~1, ad~3--4), \emph{CH} 10 (a.~3
 ad~1), \emph{CH} 5 (a.~5 ad~1); \emph{EH}
 5 (a.~2 obj.~3) and \emph{DN} 12 (a.~5 ad~2), cited by Aquinas; Augustine,
 \emph{De civitate Dei} XII.1 and XII.9 (a.~1 co., a.~4 obj.~2, a.~8 co.),
 \emph{De Trinitate} III.4.9 (a.~6 co.), \emph{De cura pro mortuis gerenda}
 16.19 (a.~8 ad~2), \emph{Enchiridion} 29 and 58;
 Gregory, \emph{Hom.\ in Ev.} 34.6--12 (a.~5 co., ad~3--4; a.~6 co.,
 obj.~4); Peter Lombard, \emph{Sent.} II d.~9 cc.~1--7 (\emph{sed contra} of
 a.~4; a.~2 obj.~2); Chrysostom, \emph{Hom.\ in Matt.} 28.3 (a.~8 obj.~3);
 Isidore, \emph{Etym.} VII.5; Bernard, \emph{De consideratione} V.4--5;
 Eph 1:20--21, Col 1:16, Isa 6:2--3, Jude 9, Matt 11:11, Judg 5:20, 1~Cor
 15:24, Heb 1:14, Matt 22:30, Luke 20:36, Rom 13:2--4, Ps 67[68]:26, Dan 10:13,
 Ps 9:5, Ps 90[91]:11, Wis 11:21; Aristotle, \emph{Metaph.} XII and
 \emph{Ethics} I (cited by Aquinas).}
 {Those who place a ``supercelestial'' hierarchy in the divine Persons (a.~1);
 the objection that the orders are by grace only (a.~4 obj.~1--3); Gregory's
 placing of the Principalities and Virtues, received by the Master under
 Dionysius's name (a.~6; \emph{Sent.} II d.~9 c.~1); those who denied that
 men can be made equal to the angels, Eustratius's tenth order beneath the
 nine, and the assignment of virgins or the perfect alone to the angelic
 orders (a.~8; \emph{In II Sent.} d.~9 q.~1 a.~8); those who took the
 demons for souls of the dead (a.~8 ad~3; Chrysostom).}

\subsection{The Ordering of the Bad Angels (\ST{I q.~109})}\label{sec:q109}

The four articles ask whether there are orders among the demons; whether
there is precedence among them; whether one enlightens another; and whether
they are subject to the precedence of the good angels (\ST{I q.~109 pr.}).
Scripture speaks of the fallen as an army with a head: ``Michael and his
angels fought with the dragon, and the dragon fought, and his angels'' (Apoc
12:7); the everlasting fire ``was prepared for the devil and his angels''
(Matt 25:41); the Pharisees name ``Beelzebub the prince of the devils,'' and
the Lord answers with the image of a kingdom, ``how then shall his kingdom
stand?'' (Matt 12:24--26).

\paragraph{There are orders among the demons (a.~1).} The \emph{sed contra}
is the Apostle: ``our wrestling is not against flesh and blood; but against
principalities and powers, against the rulers of the world of this darkness''
(Eph 6:12). Order among the angels is considered according to the grade of
nature and according to the grade of grace, and grace has two states,
imperfect, the state of merit, and perfect, the state of glory. ``If
therefore we consider the angelic orders in the light of the perfection of
glory, then the demons are not in the angelic orders, and never were. But if
we consider them in relation to imperfect grace, in that view the demons were
at the time in the orders of angels, but fell away from them,'' since all the
angels were created in grace (\ST{I q.~62 a.~3}). ``But if we consider them
in the light of nature, in that view they are still in those orders; because
they have not lost their natural gifts; as Dionysius says (Div. Nom. iv)''
(\ST{I q.~109 a.~1 co.}). The \emph{Divine Names} says it in a sentence of
great dignity: ``the whole good which was given to them was not changed, but
themselves fell from the whole good given. And the angelic gifts which were
given to them, we by no means affirm that they were changed, but they exist,
and are complete, and all luminous, although the demons themselves do not
see, through having blunted their powers of seeing good'' (\emph{DN} 4.23).

The replies preserve the goodness of the created order even in the fallen.
Aquinas cites Augustine that order belongs to the good as do mode and
species; but ``good can exist without evil; whereas evil cannot exist without
good~\dots\ so there is order in the demons, as possessing a good nature''
(ad~1).
Their ordering is not a holy principality on their side, ``because they abuse
their nature for evil''; but on the side of God who orders them ``it is
sacred; for He uses the demons for Himself'' (ad~2). The fallen are never
called Seraphim, Thrones, or Dominations, for ``the name `Seraphim' is given
from the ardor of charity; and the name `Thrones' from the Divine indwelling;
and the name `Dominations' imports a certain liberty; all of which are
opposed to sin'' (ad~3).

The third objection records that ``the demons fell from every one of the
angelic orders; as is commonly supposed'' (\latin{ut communiter dicitur}).
The Master teaches it: \latin{De ordine namque superiori Lucifer ille fuit,
quo nullus dignior conditus est. Apostolus etiam principatus et potestates
tenebrarum nominat, ostendens de ordinibus illis cecidisse} (\emph{Sent.} II
d.~9 c.~4). Gregory reads the nine precious stones that covered the prince of Tyre
(Ezek 28:13) as the nine orders, which the first angel wore as the vesture
of his adornment and outshone; the holes pierced in the stones to be bound
with gold are the capacity for love, which in him pride left unfilled, while
the others, bound together by charity, received on his fall the gift never
to be loosened from the ornament of the King (\emph{Moralia} XXXII.23.48;
\S\ref{sec:fa-gregory}).
Aquinas's commentary records that Augustine left undetermined whether the
first sinner was the highest of all, \latin{eo quod de his quae pertinent ad
Angelos, pauca voluit asserendo tradere}, but that Gregory determines
expressly that he was higher than all, \latin{cui consentit communis
sententia} (\emph{In II Sent.} d.~6 q.~1 a.~1). The first sin and the rank
of the first sinner are treated at \S\ref{sec:q63}.

\paragraph{There is precedence among the demons (a.~2).} The \emph{sed
contra} is the gloss on 1~Cor 15:24: ``While the world lasts, angels will
preside over angels, men over men, and demons over demons.'' ``Since action
follows the nature of a thing, where natures are subordinate, actions also
must be subordinate to each other,'' as the movements of lower bodies are
subject to those of the heavenly. The demons are by natural order subject
to one another; ``and hence their actions are subject to the action of those
above them, and this is what we mean by precedence~\dots\ So the very natural
disposition of the demons requires that there should be authority among
them. This agrees too with Divine wisdom, which leaves nothing inordinate,
which `reacheth from end to end mightily, and ordereth all things sweetly'
(Wisdom 8:1)'' (\ST{I q.~109 a.~2 co.}).

The objections are drawn from the nature of evil, and the replies expose it.
Their authority ``is not founded on their justice, but on the justice of God
ordering all things'' (ad~1). ``Among the proud there are always
contentions'' (Prov 13:10), yet the demons obey one another, for their
concord ``does not arise from mutual friendships, but from their common
wickedness whereby they hate men, and fight against God's justice. For it
belongs to wicked men to be joined to and subject to those whom they see to be
stronger, in order to carry out their own wickedness'' (ad~2). And rule among
them is no reward: ``That the inferior are subject to the superior, is not
for the benefit of the superior, but rather to their detriment; because since
to do evil belongs in a pre-eminent degree to unhappiness, it follows that to
preside in evil is to be more unhappy'' (ad~3).

The commentary on the \emph{Sentences} gives three grounds for the order
among the demons, their nature, the divine wisdom, and their own wickedness:
from nature some were higher than others, and sin does not take nature away;
the power of the demons to exercise men and to punish the damned is from God,
\latin{et ideo ordinata per gradus praelationis debet esse}; and because
they are enemies of mankind, \latin{amicos autem esse eos qui unius inimici
sunt, consequitur}, they are confederated to attack in concord and order. Its
\emph{sed contra} is the scales of Leviathan, of which it is written that
``One is joined to another, and not so much as any air can come between
them: They stick one to another and they hold one another fast, and shall
not be separated'' (Job 41:7--8 [41:8--9]). The replies add that the demons have no
benevolence toward one another, only concord for a deed; that some form of
right can remain in acts of injustice, \latin{sicut patet in latronibus, qui
alios injuste spoliantes, inter se aliquam formam justitiae habent, dum sibi
invicem fidem servant}; and that the higher in office bear the heavier
torment, as it is written, ``the mighty shall be mightily tormented'' (Wis
6:7) (\emph{In II Sent.} d.~6 q.~1 a.~4). Gregory names the head of this
kingdom from the last verse of the Lord's speech to Job, ``He is a king over
all the children of pride'' (Job 41:25 [41:26]): Leviathan ``smote himself
with pride alone,'' and by pride overthrew the men who followed him
(\emph{Moralia} XXXIV.23.47).

\paragraph{There is no enlightenment among the demons (a.~3).} The
\emph{sed contra} is Ecclesiasticus: ``What can be made clean by the
unclean?'' (Ecclus 34:4); enlightenment goes with cleansing and perfecting,
and cleansing does not befit the demons. ``Enlightenment properly speaking is
the manifestation of the truth in reference to God, Who enlightens every
intellect''; another manifestation of truth is speech. ``Now the demon's
perversity does not lead one to order another to God, but rather to lead away
from the Divine order; and so one demon does not enlighten another; but one
can make known his mental concept to another by way of speech'' (\ST{I
q.~109 a.~3 co.}; on speech, \S\ref{sec:q107}). The greater natural light of
the higher demons gives no ground of enlightenment, since in natural
knowledge no manifestation is needed among angels or demons: ``they know from
the first all that belongs to their natural knowledge'' (ad~2;
\S\ref{sec:q55}).

\paragraph{The good angels rule the bad (a.~4).} The \emph{sed contra}
gathers Augustine, that ``the treacherous and sinful spirit of life is ruled
by the rational, pious, and just spirit of life,'' and Gregory, that ``the
Powers are the angels to whose charge are subjected the hostile powers''
(\ST{I q.~109 a.~4 s.c.}; Augustine, \emph{De Trinitate} III.4.9; Gregory,
\emph{Hom.\ in Ev.} 34.10). ``The whole order of precedence is first and
originally in God; and it is shared by creatures accordingly as they are the
nearer to God.'' The greatest perfection, which brings a creature nearest to
God, belongs to those who enjoy God, as the holy angels do, ``of which
perfection the demons are deprived; and therefore the good angels have
precedence over the bad, and these are ruled by them'' (\ST{I q.~109 a.~4
co.}).

The replies describe the government of the fallen with judicial exactness.
``Many things concerning Divine mysteries are made known by the holy angels to
the bad angels, whenever the Divine justice requires the demons to do anything
for the punishment of the evil; or for the trial of the good; as in human
affairs the judge's assessors make known his sentence to the executioners.''
On the side of the holy angels such a revelation is an enlightenment, since
they direct it to God; on the side of the demons it is not, ``for these do not
direct it to God; but to the fulfilment of their own wickedness'' (ad~1).
``The holy angels are the ministers of the Divine wisdom. Hence as the Divine
wisdom permits some evil to be done by bad angels or men, for the sake of the
good that follows; so also the good angels do not entirely restrain the bad
from inflicting harm'' (ad~2). And an angel lower in nature presides over
demons higher in nature, ``because the power of Divine justice to which the
good angels cleave, is stronger than the natural power of the angels,'' as
among men ``the spiritual man judgeth all things'' (1~Cor 2:15) (ad~3).
Bernard gives the same office its pastoral face: ``Let us suppose that the
Powers are their superiors, and that by their vigour the power of darkness is
checked, and the malignant spirits of this lower air are restrained, that
they may not do harm to their full intent: that they may not be able to show
their malignity, except for beneficial ends'' (\emph{De consideratione}
V.4.8). ``And that great dragon was cast out, that old serpent, who is called
the devil and Satan~\dots\ and his angels were thrown down with him'' (Apoc
12:9). The assaults of the demons and the defence of men are treated at
\S\ref{sec:assaults}.

\angeldossier{\ST{I q.~109 aa.~1--4}; parallels \emph{In II Sent.} d.~6
q.~1 aa.~1,~4.}
 {Scripture teaches that the fallen angels have a head and form a kingdom
 (Matt 12:26; Matt 25:41; Apoc 12:7--9; Eph 6:12). Common teaching that
 there is order and precedence among the demons, founded on their unlost
 natural gifts and on the divine justice (aa.~1--2; \emph{DN} 4.23; the
 gloss on 1~Cor 15:24), and that the good angels rule the bad (a.~4;
 Augustine, \emph{De Trin.} III.4.9; Gregory, \emph{Hom.\ in Ev.}\ 34.10). Common
 teaching that some fell from every order and that the first sinner was the
 highest of all (a.~1 obj.~3; Gregory, \emph{Moralia} XXXII.23.47--48;
 Lombard, \emph{Sent.} II d.~9 c.~4; \emph{In II Sent.} d.~6 q.~1 a.~1,
 recording Augustine's reserve). Thomist position that the demons remain in
 the orders by nature but never were in them by glory (a.~1), that their
 ordering is sacred on God's side only (a.~1 ad~2), and that among them there
 is speech but no enlightenment (a.~3).}
 {Eph 6:12 (\emph{sed contra} of a.~1); the gloss on 1~Cor 15:24
 (\emph{sed contra} of a.~2); Ecclus 34:4 (\emph{sed contra} of a.~3);
 Augustine, \emph{De Trinitate} III.4.9, and Gregory,
 \emph{Hom.\ in Ev.} 34.10 (\emph{sed contra} of a.~4); Dionysius,
 \emph{DN} 4.23 (a.~1 co.); Augustine,
 \emph{De natura boni} (a.~1 obj.~1, cited by Aquinas); Gregory,
 \emph{Moralia} XXXII.23.47--48 and XXXIV.23.47; Peter Lombard, \emph{Sent.} II d.~9
 c.~4; Bernard, \emph{De consideratione} V.4.8; Prov 13:10, Wis 8:1, 1~Cor
 2:15, Job 41:7--8 [41:8--9], 41:25 [41:26], Wis 6:7, Ezek 28:13, Matt 12:24--26, Matt 25:41,
 Apoc 12:7--9; Aristotle, \emph{Ethics} III and X (a.~4 ad~3, cited by
 Aquinas).}
 {The objections that order belongs to the good alone and that there can be
 no hierarchy where there is no holiness (a.~1 obj.~1--2); that there is no
 precedence without concord and justice (a.~2 obj.~1--3); that the higher
 demons, having natural light in abundance, enlighten the lower (a.~3
 obj.~1--2); that negligence would be charged to good angels presiding over
 the bad (a.~4 obj.~2). Augustine, as Aquinas reports him, leaves
 undetermined whether the first sinner was the highest of all (\emph{In II
 Sent.} d.~6 q.~1 a.~1).}
